% Copyright (c) 2026 by C. Klukas.
% Licensed under the MIT License. See LICENSE for details.

A \paglet{} is a WebAssembly module together with one instance of it.
This chapter takes that sentence apart: what names a \paglet{} and who
runs it, how it receives and sends messages, how payloads are typed, how
it reaches other \paglet{}s, and what it can and cannot do. The chapters
that follow build on these definitions: the lifecycle
(Section~\ref{sec:lifecycle}), moves (Section~\ref{sec:movement}),
communication while a \paglet{} moves (Section~\ref{sec:communication})
and the mesh (Section~\ref{sec:mesh}).

\subsection{A module and one instance}\label{sec:model-instance}

The \emph{module} is a core WebAssembly module for \code{wasm32}, compiled
by clang from C++ against the guest SDK (\code{cpp/sdk/}). It has exactly
one linear memory, exported as \code{memory}, and a small set of exports
the host calls (Section~\ref{sec:model-imports}). The host names a module
by the SHA-256 hash of its bytes, written in hex; adding the same bytes
twice is a no-op, and every \paglet{} record, every sender record and
every passport refers to its module by this hash.

The \emph{instance} is one WebAssembly instance of that module. The SDK
keeps a single static object of the \paglet{}'s class in the module's
linear memory. That object, everything it owns on the heap, the
registered handlers and the continuations of outstanding requests are all
bytes of that memory. The \paglet{}'s state is therefore its linear
memory plus the module's exported mutable globals (for C++ guests, the
stack pointer). A \emph{memory image} of these is all a host needs to
store a \paglet{}, to copy it, or to send it to another host
(Section~\ref{sec:movement}); no code in a \paglet{} serializes its
state.

Three values identify a \paglet{} to its host and to the mesh:

\begin{itemize}
\item The \emph{\paglet{} ID}: 32 lowercase hex digits, 16 random bytes
  drawn when the \paglet{} is created, or chosen in advance when a
  passport names it.
\item The \emph{owner}: a string the host assigns. On a single host it is
  \code{local}; in a mesh it is the key ID of the owner whose key signed
  the launch. Children and clones have the owner of their creator.
\item The \emph{trust class}: \code{roaming}, \code{resident} or
  \code{system} (Table~\ref{tab:model-trust}). The host assigns it; a
  \paglet{} never claims its own class, and a \paglet{} that arrives from
  another host is always roaming. Children are roaming. A clone keeps its
  creator's class, except that a clone of a system \paglet{} is roaming;
  a clone made on another host (Section~\ref{sec:movement}) is only
  allowed for a roaming \paglet{}.
\end{itemize}

\begin{table}[htbp]
\centering\small
\begin{tabular}{lp{68mm}lp{32mm}}
\toprule
Class & Origin & Moves & Import set \\
\midrule
\code{system} & \raggedright built into the host (native C++), or a module the ledger trusts for this class & never & \raggedright standard plus \code{paglets\_sys} \tabularnewline
\code{resident} & \raggedright meant for modules an admin installs on one host; the ledger can admit modules to the class, but no command or mesh operation creates a resident \paglet{} today & no & standard \tabularnewline
\code{roaming} & \raggedright created by an owner or by another \paglet{}, or arrived from another host & yes & standard \tabularnewline
\bottomrule
\end{tabular}
\caption{Trust classes. In the current release every system \paglet{} is native; the runtime's C++ API (\code{CreateOptions::trust}) accepts every class and the tests use it, but no command or mesh operation creates a resident \paglet{}. The classes exist in the runtime, the ledger's module trust and the sender records.}
\label{tab:model-trust}
\end{table}

A \paglet{} launched into a mesh carries a \emph{passport}: a document
signed by the owner's key that names the mesh, the owner, the module hash
and the \paglet{} ID, extended by a host-signed link for every child or
clone. A host verifies it against the ledger before it accepts the
\paglet{}, and it travels with the \paglet{}. A \paglet{} without one,
such as one started by \code{paglets-host run} on a single host, cannot
move. Section~\ref{sec:mesh-passports} gives the contents and the checks.

Figure~\ref{fig:model-architecture} shows where an instance runs. The host
process, \code{paglets-host}, starts one \emph{worker process} per
scheduler lane (two by default) and places each active \paglet{} on one of
them. The operating system sandboxes the workers: a seccomp filter on
Linux, a sandbox profile on macOS, a job object with a restricted token on
Windows. A worker can neither open files nor create sockets nor start
programs (on Windows the restricted token stops new processes and writes
to ordinary files, but reading files and network access are not
restricted yet); every effect of a \paglet{} passes through its host.
Calls between the two processes go over two shared-memory rings, one per
direction; a socket pair (two pipes on Windows) carries wake-ups and
terminations. The embedded WAMR engine runs the instances. For tests and debugging,
\code{--in-process} runs them inside the host process instead.

\begin{figure}[htbp]
\centering
\pgfdeclarelayer{background}
\pgfsetlayers{background,main}
\begin{tikzpicture}[node distance=3.5mm and 36mm]
  \node[plain,text width=44mm] (module) {\textbf{Paglet module} (Wasm)\\[1pt]\footnotesize linear memory: the \paglet{} object, its handlers, its continuations};
  \node[plain,text width=44mm,below=of module] (sdk) {\textbf{Guest SDK}\\[1pt]\footnotesize \code{Paglet}, \code{Router}, \code{Message}, \code{Endpoint}, MessagePack codecs};
  \node[plain,text width=44mm,below=of sdk] (wamr) {\textbf{WAMR engine}\\[1pt]\footnotesize instances, traps, time budgets, memory images};
  \begin{pgfonlayer}{background}
  \node[box,fit=(module)(sdk)(wamr),inner sep=7pt,label={[font=\small\bfseries\color{pagletsBlue}]above:Worker process (sandboxed)}] (worker) {};
  \end{pgfonlayer}

  \node[plain,text width=44mm,right=of module] (runtime) {\textbf{Runtime}\\[1pt]\footnotesize mailboxes, capability tables, timers, images, scheduler};
  \node[plain,text width=44mm,below=of runtime] (system) {\textbf{System \paglet{}s}\\[1pt]\footnotesize files, directory, grants, mesh-info, \ldots};
  \node[plain,text width=44mm,below=of system] (node) {\textbf{Mesh node}\\[1pt]\footnotesize moves, location records, ledger, channels};
  \begin{pgfonlayer}{background}
  \node[box,fit=(runtime)(system)(node),inner sep=7pt,label={[font=\small\bfseries\color{pagletsBlue}]above:Host process (\code{paglets-host})}] (host) {};
  \end{pgfonlayer}

  \draw[arrow] (host.west|-module) -- node[event,above,text width=28mm,font=\scriptsize\color{pagletsMuted}]{\code{paglets\_on\_event}\\\code{paglets\_on\_message}} (worker.east|-module);
  \draw[arrow] (worker.east|-sdk) -- node[event,above,text width=28mm,font=\scriptsize\color{pagletsMuted}]{\code{send}, \code{request}, \code{reply}, \code{lifecycle}, \ldots} (host.west|-sdk);
  \draw[both] (worker.east|-wamr) -- node[event,above,text width=28mm,font=\scriptsize\color{pagletsMuted}]{memory images} (host.west|-wamr);
  \node[event,text width=60mm] at ($(worker.south east)!0.5!(host.south west)+(0,-6mm)$) {Paglet ABI v1: MessagePack documents in linear memory, carried between the processes over shared-memory rings};
  \draw[dashedarrow] (host.east) -- ++(7mm,0) node[event,right,text width=12mm]{other hosts};
\end{tikzpicture}
\caption{A \paglet{} instance in a worker process and the host runtime. The host calls the module's handler exports; the SDK calls the fourteen imports of the \code{paglets} module, and the guest's C library the WASI subset of Table~\ref{tab:model-imports}. Nothing else crosses the boundary.}
\label{fig:model-architecture}
\end{figure}

\subsection{A first \paglet{}}\label{sec:model-hello}

Listing~\ref{lst:hello} is the smallest example of the repository. A
\paglet{} is a class derived from \code{paglets::Paglet}, registered with
\code{PAGLETS\_PAGLET} at namespace scope. Its constructor runs once,
when the host delivers the \code{created} event, and never again: after a
restore from a memory image the constructed object is already in the
memory. The constructor is where handlers are registered.

% snippet: cpp/examples/hello/hello.cpp [float,caption={The hello \paglet{}},label={lst:hello}]
\begin{lstlisting}[style=paglets-cpp,float,caption={The hello \paglet{} \src{cpp/examples/hello/hello.cpp}},label={lst:hello}]
// Smallest paglet: answers "greet" requests and counts them.

#include <paglets/paglet.hpp>

#include <string>

namespace {

class Hello : public paglets::Paglet {
public:
    Hello() {
        router().on<std::string>("greet", [this](const std::string& name, paglets::Message& m) {
            ++greetings_;
            m.reply("hello, " + name + " (greeting " + std::to_string(greetings_) + ")");
        });
    }

    void on_created(paglets::Started&) override {
        auto info = paglets::self_info();
        paglets::log("hello paglet created as " + (info ? info->id : std::string("?")));
    }

private:
    std::uint32_t greetings_ = 0;
};

}  // namespace

PAGLETS_PAGLET(Hello)
\end{lstlisting}
% /snippet

\code{router().on<std::string>("greet", ...)} registers a handler for
requests named \code{greet} whose body decodes as a MessagePack string.
The handler receives the decoded body and the \code{Message}, and answers
with \code{m.reply(...)}, which the SDK encodes as MessagePack again. The
count of greetings is an ordinary member: nothing saves it, it is part of
the memory image. \code{on\_created} is one of the lifecycle hooks of
Section~\ref{sec:lifecycle}; here it asks for the \paglet{}'s own record
with \code{self\_info()} and writes the ID to the host log, where a
\paglet{}'s \code{stdout} and \code{stderr} go as well. One CMake call
builds the module: \code{paglets\_add\_module(hello SOURCES hello/hello.cpp)}
(the repository's line adds \code{INTERFACE hello/interface.json}, the
interface description of Section~\ref{sec:tools}).
\code{paglets-host run hello.wasm} with \code{--call greet '"world"'}
creates a \paglet{} from it, sends the request and prints the reply as
JSON: \code{greet -> "hello, world (greeting 1)"}. Bodies on the command
line are JSON; the host converts them to MessagePack and the replies back.

\subsection{Messages, requests and timers}\label{sec:model-messages}

Every \paglet{} has a mailbox. The host delivers its messages one at a
time, by priority (0 to 7, highest first, default 3) and in arrival order
within a priority; a mailbox holds at most 1024 messages. A \paglet{}
never sees two calls at once: handlers and lifecycle hooks run one after
the other, never concurrently, and the host takes memory images only
between two calls. A \paglet{} shares memory with nothing; its only
effects are the messages it sends, the \paglet{}s it creates, the timers
it sets, its log lines and the lifecycle operations it requests. A handler
must not block, because nothing it could wait for arrives while it runs: a
reply, a service's answer or a timer is a later message. A handler that
runs longer than its time budget (5 seconds of elapsed time by default) is
terminated and the \paglet{} fails.

A delivered \code{Message} has a kind (\code{message}, \code{request},
\code{reply} or \code{timer}), a name of up to 128 bytes, a payload, a
priority, the capabilities that came with it, and for messages from other
\paglet{}s, or from the host's own tools (ID \code{host}, trust
\code{system}), a \emph{sender record} the host stamps on it: the
sender's ID, owner, trust class, module hash and host. Replies, timers and
reports from the host carry none. Three things can be sent:

\begin{itemize}
\item A \emph{message} is one-way: \code{endpoint.send(name, body)}. The
  call fails at once with \code{denied}, \code{expired} or \code{quota}
  when the endpoint does not allow the message, with \code{not\_found}
  when the target has ended, and with \code{quota} when its mailbox is
  full. With \code{SendOptions::async} the call returns as soon as the
  document is well-formed and a failure arrives later as the hook
  \code{on\_undelivered}; in a worker process this saves the round trip to
  the host.
\item A \emph{request} expects exactly one reply.
  \code{request(name, body, on\_reply)} on an endpoint returns a
  correlation ID and keeps \code{on\_reply}, the continuation. The
  receiver gets the message with a one-shot \emph{reply capability}. The
  reply arrives at priority 7 with a status: 0 and the payload for a real
  reply; otherwise \code{timeout} (30 seconds by default),
  \code{unknown\_message} (the receiver had no handler), \code{failed} (the
  receiver trapped), \code{gone} (the receiver ended or dropped the reply
  capability), or the error code the receiver's handler returned. A
  request is never asynchronous in the sense above: the call waits for the
  host to accept it.
\item A \emph{timer}: \code{after(delay\_ms, name, body)} delivers a
  message of kind \code{timer} to the \paglet{} itself after the delay.
  The returned capability cancels the timer when dropped. A \paglet{}
  holds at most 64 pending timers; they survive deactivation and wake a
  deactivated \paglet{}.
\end{itemize}

The continuation is a \code{std::function} that the SDK stores in a map
keyed by the correlation ID. The map lives in linear memory, so an
outstanding request survives deactivation and moves: the reply reaches its
continuation wherever the \paglet{} is when it arrives.
Listing~\ref{lst:pingpong-bounce} shows the shape. \code{pong\_} is the
parent's endpoint to the child, a capability
(Section~\ref{sec:model-capabilities}); \code{pending\_run\_} holds the
reply capability of the \code{run} request that started the rounds
(Listing~\ref{lst:pingpong-roles}). The parent of the ping-pong example
sends \code{ping} to its child and names the function that runs with the
reply; that function checks the status, decodes the body and sends the
next \code{ping}. No handler waits.

% snippet: cpp/examples/ping_pong/ping_pong.cpp#bounce [float,caption={A request with a continuation, and a deferred reply answered from it},label={lst:pingpong-bounce}]
\begin{lstlisting}[style=paglets-cpp,float,caption={A request with a continuation, and a deferred reply answered from it \src{cpp/examples/ping\_pong/ping\_pong.cpp}},label={lst:pingpong-bounce}]
void bounce(std::uint32_t n) {
    if (n >= target_) {
        (void)paglets::reply_to(pending_run_, n);
        (void)pong_.send("stop");
        return;
    }
    auto sent = pong_.request("ping", n, [this](paglets::Message& reply) {
        std::uint32_t value = 0;
        if (!reply.ok() || !reply.decode(value)) {
            paglets::log(paglets::abi::LogLevel::error,
                         "ping failed: " + std::string(paglets::abi::error_name(reply.status())));
            (void)pending_run_.drop();
            return;
        }
        bounce(value);
    });
    if (!sent) (void)pending_run_.drop();
}
\end{lstlisting}
% /snippet

A request is answered once. The handler answers it directly with
\code{reply} or \code{reply\_async} on the message, or it keeps the reply
capability for later: \code{defer\_reply} returns it, and
\code{reply\_to} answers with it from any later handler. This
is the \emph{deferred reply} of Listing~\ref{lst:pingpong-roles}: the
\code{run} handler cannot know its answer, which exists only after many
rounds, so it defers the reply, and \code{bounce} answers it when the
target is reached. A reply capability is an ordinary capability: it can be
transferred to another \paglet{}, which then answers in the original
receiver's place. Dropping it, as the error path of
Listing~\ref{lst:pingpong-bounce} does, answers the requester \code{gone}
at once instead of letting it wait for its timeout. A request that a
handler handled but neither answered nor deferred is answered \code{gone}
as well: the SDK drops the reply capability when the handler returns.

% snippet: cpp/examples/ping_pong/ping_pong.cpp#roles [float,caption={The ping-pong handlers: a child from the same module, a deferred reply, an asynchronous reply},label={lst:pingpong-roles}]
\begin{lstlisting}[style=paglets-cpp,float,caption={The ping-pong handlers: a child from the same module, a deferred reply, an asynchronous reply \src{cpp/examples/ping\_pong/ping\_pong.cpp}},label={lst:pingpong-roles}]
PingPong() {
    // Parent role.
    router().on<std::uint32_t>("run", [this](std::uint32_t rounds, paglets::Message& m) -> std::int32_t {
        if (pending_run_) return paglets::abi::bad_state;
        auto child = paglets::create_child(std::uint32_t{0});
        if (!child) return child.error();
        pong_ = *child;
        target_ = rounds;
        pending_run_ = m.defer_reply();
        bounce(0);
        return paglets::abi::handled;
    });
    // Child role. The reply is asynchronous: in a worker process it does
    // not wait for the host, and a failure would reach on_undelivered.
    router().on<std::uint32_t>("ping", [](std::uint32_t n, paglets::Message& m) { m.reply_async(n + 1); });
    router().on("stop", [](paglets::Message&) { (void)paglets::dispose(); });
}
\end{lstlisting}
% /snippet

\subsection{The router}\label{sec:model-router}

\code{Paglet::on\_message} receives every delivery except replies that
have a continuation and the host's reports (\code{on\_undelivered},
\code{on\_move\_failed}, \code{on\_recovered}). Its default passes the
message to the \paglet{}'s \code{Router}, a map from message names to
handlers. \code{router().on(name, f)} registers \code{f} for a name;
\code{router().otherwise(f)} catches the rest. Without either, a message
is not handled: a one-way message is dropped, a request is answered
\code{unknown\_message}. Timer messages go through
the router under the name given to \code{after}. A handler either returns
\code{void}, which counts as handled, or an \code{std::int32\_t}:
\code{abi::handled} (0), \code{abi::not\_handled} (1), or a negative error
code. For a request the result decides the reply when the handler has not
replied itself: \code{not\_handled} makes the host answer
\code{unknown\_message}, an error code is answered as that status. The
typed form of Listing~\ref{lst:router-typed} decodes the payload into a
\code{T} before the handler runs and answers \code{malformed} when the
body does not decode as a complete value of that type; its handler takes
\code{const T\&} and the \code{Message}.

% snippet: cpp/sdk/include/paglets/paglet.hpp#router-typed [float,caption={The typed handler registration of the router},label={lst:router-typed}]
\begin{lstlisting}[style=paglets-cpp,float,caption={The typed handler registration of the router \src{cpp/sdk/include/paglets/paglet.hpp}},label={lst:router-typed}]
template <class T, class F>
void on(std::string name, F&& f) {
    handlers_[std::move(name)] = [g = std::forward<F>(f)](Message& m) -> std::int32_t {
        T body{};
        if (!m.decode(body)) return abi::malformed;
        if constexpr (std::is_same_v<std::invoke_result_t<decltype(g)&, const T&, Message&>, void>) {
            g(body, m);
            return abi::handled;
        } else {
            return g(body, m);
        }
    };
}
\end{lstlisting}
% /snippet

\subsection{Typed payloads: schema headers and contracts}\label{sec:model-schemas}

Payloads are MessagePack. The SDK encodes and decodes \code{bool}, the
integer types, \code{float}, \code{double}, \code{std::string},
\code{Bytes} (as MessagePack binary), \code{std::vector<T>} and
\code{std::optional<T>} (nil when empty) by itself, which is why the hello
and ping-pong \paglet{}s need no message definitions. Structured messages
are plain C++ structs and enums in a \emph{schema header}: a struct is
encoded as a map with its field names as keys, an enum as the name of its
enumerator, and decoders skip keys they do not know, so a field added
later keeps old senders and receivers compatible.
Listing~\ref{lst:counter-messages} is the schema header of the counter
example; Listing~\ref{lst:counter-handlers} shows its handlers using the
types.

% snippet: cpp/examples/counter/messages.hpp [float,caption={The counter's schema header: wire types as plain structs and enums},label={lst:counter-messages}]
\begin{lstlisting}[style=paglets-cpp,float,caption={The counter's schema header: wire types as plain structs and enums \src{cpp/examples/counter/messages.hpp}},label={lst:counter-messages}]
// Message types of the counter example paglet. Plain structs: the host encodes
// them through C++26 reflection, the guest through code generated from this
// header by schema_gen.

#pragma once

#include <cstdint>
#include <optional>
#include <string>
#include <vector>

namespace counter_msgs {

enum class Mode { add, multiply };

struct Increment {
    std::int64_t by = 1;
    std::string note;
    Mode mode = Mode::add;
};

struct Status {
    std::int64_t value = 0;
    std::uint32_t history_size = 0;
    std::vector<std::string> recent_notes;
    std::optional<double> average_step;
    std::uint32_t memory_pages = 0;
};

// Exercises every supported field type; the guest echoes it unchanged.
struct Echo {
    bool flag = false;
    std::int8_t tiny = 0;
    std::int32_t medium = 0;
    std::uint64_t big = 0;
    float ratio = 0.0f;
    double precise = 0.0;
    std::string text;
    std::vector<std::uint8_t> blob;
    std::vector<Increment> items;
    std::optional<std::string> label;
    std::optional<Status> nested;
    Mode mode = Mode::add;
};

// Start arguments (optional): the value to begin with.
struct Start {
    std::int64_t value = 0;
};

// Allocates and keeps the given amount of memory (image size experiments).
struct Bloat {
    std::uint32_t kilobytes = 0;
    std::uint8_t fill = 0;
};

}  // namespace counter_msgs
\end{lstlisting}
% /snippet

% snippet: cpp/examples/counter/counter.cpp#handlers [float,caption={The counter's handlers},label={lst:counter-handlers}]
\begin{lstlisting}[style=paglets-cpp,float,caption={The counter's handlers \src{cpp/examples/counter/counter.cpp}},label={lst:counter-handlers}]
Counter() {
    router().on<Increment>("increment", [this](const Increment& inc, paglets::Message& m) {
        value_ = inc.mode == Mode::add ? value_ + inc.by : value_ * inc.by;
        history_.push_back(inc);
        m.reply(status());
    });
    router().on("status", [this](paglets::Message& m) { m.reply(status()); });
    router().on<Echo>("echo", [](const Echo& echo, paglets::Message& m) { m.reply(echo); });
    router().on<Bloat>("bloat", [this](const Bloat& bloat, paglets::Message& m) {
        ballast_.emplace_back(std::size_t{bloat.kilobytes} * 1024, bloat.fill);
        m.reply(status());
    });
    router().on("schema", [](paglets::Message& m) { m.reply(std::string(paglets_schema_json)); });
    router().on<std::string>("log", [this](const std::string& text, paglets::Message& m) {
        paglets::log(text);
        m.reply(status());
    });
}
\end{lstlisting}
% /snippet

The two sides of the wire get their codecs differently. The host is
compiled with GCC~16 and C++26 reflection (\code{-freflection}); its codec
(\code{wire/reflect.hpp}) walks the members of a struct at compile time
and needs no code per type. The guest is compiled by clang for
\code{wasm32-wasip1}, which has no reflection. So the build compiles the
\emph{schema generator} (\code{tools/schema\_gen}) with GCC~16 once per
schema header, named in CMake by \code{SCHEMA\_HEADER} and
\code{SCHEMA\_NAMESPACE}; it reflects over the namespace and writes
\code{<module>.schema.gen.hpp}, plain C++ encoders and decoders for every
type, which the guest includes. The two codecs agree byte for byte. The
generator also writes a JSON descriptor of the types, as
\code{paglets\_schema\_json} in the header and as a file for tools. A host
built without reflection (the \code{linux-host-only} preset) cannot run
the generator and skips such modules with a warning.

A \emph{contract} names a service's operations with their request and
reply types. Listing~\ref{lst:calc-contract} declares one: a struct with a
\code{static constexpr std::string\_view service} and one member function
per operation, each taking the request type and returning the reply type.
The functions are declared, never defined; they are a description. The
system services of Section~\ref{sec:mesh} are contracts of this form in
\code{cpp/common/include/paglets/services/}.

% snippet: cpp/tests/guests/calc_contract.hpp [float,caption={A service contract},label={lst:calc-contract}]
\begin{lstlisting}[style=paglets-cpp,float,caption={A service contract \src{cpp/tests/guests/calc\_contract.hpp}},label={lst:calc-contract}]
// A test service contract: served natively by the host and by the calc_user
// guest, called by both (tests/test_contract.cpp).

#pragma once

#include <cstdint>
#include <string>
#include <string_view>
#include <vector>

namespace calc {

struct AddRequest {
    std::int64_t a = 0;
    std::int64_t b = 0;
};

struct AddReply {
    std::int64_t sum = 0;
};

struct DivRequest {
    std::int64_t a = 0;
    std::int64_t b = 0;
};

struct DivReply {
    std::int64_t quotient = 0;
    std::string served_by;
};

struct Contract {
    static constexpr std::string_view service = "calc";
    AddReply add(const AddRequest&);
    // invalid_argument for b == 0.
    DivReply divide(const DivRequest&);
};

}  // namespace calc
\end{lstlisting}
% /snippet

For a contract the generator emits a \code{Client}, whose operations take
a continuation that receives a decoded \code{Result<Reply>} together with
the reply message (Listing~\ref{lst:calc-client}), and
\code{serve(router, impl)}, which registers one handler per operation
around ordinary member functions (Listing~\ref{lst:calc-operations}). The
same header serves natively on the host;
Section~\ref{sec:mesh-contracts} shows that side.

% snippet: cpp/tests/guests/calc_user.cpp#operations [float,caption={Serving the contract: the operations the generated \code{serve} registers},label={lst:calc-operations}]
\begin{lstlisting}[style=paglets-cpp,float,caption={Serving the contract: the operations the generated \code{serve} registers \src{cpp/tests/guests/calc\_user.cpp}},label={lst:calc-operations}]
paglets::Result<calc::AddReply> add(const calc::AddRequest& q, Message&) { return calc::AddReply{q.a + q.b}; }

paglets::Result<calc::DivReply> divide(const calc::DivRequest& q, Message&) {
    if (q.b == 0) return std::unexpected(paglets::abi::invalid_argument);
    return calc::DivReply{q.a / q.b, "guest"};
}
\end{lstlisting}
% /snippet

% snippet: cpp/tests/guests/calc_user.cpp#client [float,caption={Calling the contract through the generated client},label={lst:calc-client}]
\begin{lstlisting}[style=paglets-cpp,float,caption={Calling the contract through the generated client \src{cpp/tests/guests/calc\_user.cpp}},label={lst:calc-client}]
router().on<calc::DivRequest>("use_service", [](const calc::DivRequest& q, Message& m) -> std::int32_t {
    auto endpoint = paglets::service("calc");
    if (!endpoint) return endpoint.error();
    paglets::Capability pending = m.defer_reply();
    calc::Client client{*endpoint};
    auto sent = client.divide(q, [pending](paglets::Result<calc::DivReply> r, Message&) mutable {
        if (r) {
            paglets::reply_to(pending, *r);
        } else {
            paglets::reply_to(pending, "error:" + std::string(paglets::abi::error_name(r.error())));
        }
    });
    if (!sent) return sent.error();
    return paglets::abi::handled;
});
\end{lstlisting}
% /snippet

\subsection{Capabilities and endpoints}\label{sec:model-capabilities}

A \paglet{} has no address. Its ID names it in records and logs, but no
import takes an ID as a destination. The only way to reach another
\paglet{}, a service or a resource is a \emph{capability}: a positive
\code{int32} handle into a table the host keeps for each \paglet{}. The
guest cannot forge a handle, since the host checks every use against the
table, and the table is host state that travels with the \paglet{} in its
record, not inside the memory image. A capability has a kind
(\code{endpoint}, \code{reply}, \code{timer}, and the resource kinds
\code{dir}, \code{artifact}, \code{topic} and \code{pin} that system
\paglet{}s hand out; the ABI also names a \code{file} kind, which no
service of the current release hands out), a target, a list of operations, a
transferable flag, and optionally an expiry, a number of uses left and a
badge. \code{inspect()} returns this description.

An \emph{endpoint} is a capability to send to one \paglet{} or service;
its operations are the message names it admits, or \code{*}. A \paglet{}
gets endpoints in a few ways only: handle 1 is always its endpoint to
itself; \code{create\_child} and \code{clone} return an endpoint to the new
\paglet{}; a message, a reply or the \code{Started} event of a child or
clone can carry capabilities (\code{Message::take\_cap},
\code{Started::caps}); every new \paglet{} receives non-transferable
endpoints to the host's system services that offer one, listed in
\code{self\_info().services} and found with \code{paglets::service(name)};
and the \code{directory} service answers lookups of published names with
endpoints. Passing a capability in a message transfers it: the sender's
handle is invalid once the call succeeds and the receiver gets a handle of
its own. Besides handle 1 and the service endpoints, a child gets exactly
the capabilities its creator passes; it has no endpoint to its creator
unless one is among them. \code{derive} makes a weaker copy: fewer operations, an
earlier expiry, fewer uses, no transfer right, a badge the receiver sees
on every message, or for \code{dir} capabilities a path below the
directory; rights only shrink. \code{drop} releases a capability; dropping
a timer cancels it, dropping a reply capability answers \code{gone}.
Capabilities shown to a system \paglet{} for one message are \emph{lent}
(\code{SendOptions::lend}) and stay with the sender. A \paglet{} holds at
most 1024 capabilities. How endpoints keep working while their target
moves is the subject of Section~\ref{sec:communication}.

\subsection{What a \paglet{} can and cannot do}\label{sec:model-imports}

The host checks a module statically before it instantiates it. The
version marker \code{paglets\_abi\_v1} must be exported, and so must
\code{memory}, the allocator pair \code{paglets\_alloc} and
\code{paglets\_free}, and the two handlers \code{paglets\_on\_event} and
\code{paglets\_on\_message}; every mutable global must be exported so that
images can capture it; and every import must be in the import set of the
trust class (Table~\ref{tab:model-imports}). A module that imports
anything else is rejected at load time with a message naming the import.
The import set is the whole of what a \paglet{} can do. There is no file
system: the WASI subset preopens no directories, and \code{fd\_write} on
\code{stdout} and \code{stderr} goes to the host log. There are no
sockets. \code{proc\_exit} fails the \paglet{}. Host resources are reached
by message, through the system \paglet{}s of Section~\ref{sec:mesh},
which hand out resource capabilities under the mesh policy. The module
\code{paglets\_sys}, meant for system \paglet{}s written as Wasm modules,
is reserved: in the current release every system \paglet{} is native.

\begin{table}[htbp]
\centering\small
\begin{tabular}{lp{95mm}}
\toprule
Module & Imports \\
\midrule
\code{paglets} & \raggedright \code{send}, \code{request}, \code{reply}; \code{cap\_derive}, \code{cap\_drop}, \code{cap\_inspect}, \code{cap\_list}; \code{create\_child}, \code{lifecycle} (dispose, deactivate, clone, dispatch, checkpoint, restart), \code{timer\_set}; \code{log}, \code{self\_info}, \code{now}, \code{random} \tabularnewline
\code{wasi\_snapshot\_preview1} & \raggedright \code{args\_get}, \code{args\_sizes\_get}, \code{environ\_get}, \code{environ\_sizes\_get} (all empty); \code{clock\_res\_get}, \code{clock\_time\_get}, \code{random\_get}, \code{sched\_yield}; \code{proc\_exit}; \code{fd\_write}, \code{fd\_close}, \code{fd\_seek}, \code{fd\_fdstat\_get}, \code{fd\_prestat\_get}, \code{fd\_prestat\_dir\_name} \tabularnewline
\code{paglets\_sys} & \raggedright reserved for system \paglet{}s as Wasm modules; a roaming or resident module that imports from it is rejected \tabularnewline
\bottomrule
\end{tabular}
\caption{The import set of roaming and resident \paglet{}s: the fourteen functions of the \code{paglets} module and a WASI subset without files or sockets.}
\label{tab:model-imports}
\end{table}

Everything a \paglet{} sends or receives is a MessagePack document in its
own linear memory: the host copies a document into a buffer the guest
allocates, and an import passes a pointer and a length that the host
checks against the memory bounds. Documents are limited to 1~MB;
Table~\ref{tab:model-limits} lists the other default limits. A trap, such
as an out-of-bounds access, \code{unreachable}, stack exhaustion or an
exceeded time budget, fails the \paglet{}: the host discards the instance,
answers every request waiting on it with \code{failed}, and ends it. A
trap is the \paglet{}'s own fault, so it does not resume from a stored
image. When a worker process ends for another reason, the \paglet{} in
the call fails as well; the other \paglet{}s on that worker resume from
their last image, or fail if they have none.

\begin{table}[htbp]
\centering\small
\begin{tabular}{lr}
\toprule
Limit & Default \\
\midrule
Message document & 1 MB \\
Linear memory per \paglet{} & 16 MB \\
Time per handler call (elapsed) & 5 s \\
Mailbox length & 1024 messages \\
Capabilities per \paglet{} & 1024 \\
Pending timers per \paglet{} & 64 \\
Children and clones per handler call & 16 \\
Paglets of one owner per host & 10\,000 \\
\bottomrule
\end{tabular}
\caption{Default limits of a host (\code{runtime::Config}, \code{wasm::Limits} and \code{abi::max\_document\_size}).}
\label{tab:model-limits}
\end{table}

These rules give the model its shape. A \paglet{} is a single-threaded
object that reacts to messages, holds its whole state in one memory, and
reaches the outside only through handles the host gave it. The host can
stop it between any two messages, write its memory to disk, and start it
again elsewhere, and the \paglet{} cannot tell the difference except
through the events of the next section.
